% ============================================================
\chapter{Curvature --- Riemann Tensor}
\label{ch:curvature}
% ============================================================

Curvature is the quantity that distinguishes one geometry from another. On a plane, parallel lines remain parallel; on a sphere, they eventually meet; on a saddle surface, they diverge. Riemann understood, as early as 1854, that this deviation of geodesics is measured by a tensorial object---the \emph{curvature tensor}---which encodes, at each point and in each tangent plane, how space curves. This tensor, formalised by Christoffel and Ricci-Curbastro, is the central ingredient of Einstein's general relativity: the curvature of spacetime \emph{is} gravitation.

\section{The Riemann curvature tensor}

\begin{definition}[Curvature tensor]
The \emph{Riemann curvature tensor} of the Levi-Civita connection $\nabla$ is
the map $R : \mathfrak{X}(M)^3 \to \mathfrak{X}(M)$ defined by:
\[
    R(X,Y)Z = \nabla_X \nabla_Y Z - \nabla_Y \nabla_X Z - \nabla_{[X,Y]} Z.
\]
\end{definition}

\begin{intuition}[title=\textbf{Curvature measures non-commutativity}]
The Riemann tensor measures the failure of covariant differentiation to commute:
in flat space, $\nabla_X \nabla_Y = \nabla_Y \nabla_X$ (modulo the Lie bracket).
Curvature is the obstruction to parallel transport being path-independent.
\end{intuition}

\begin{proposition}
$R$ is a $(1,3)$-tensor: it is $C^\infty(M)$-multilinear in $X$, $Y$, and $Z$.
\end{proposition}

\begin{proof}
We check $C^\infty(M)$-linearity in $Z$ (the others are similar):
\begin{align*}
    R(X,Y)(fZ) &= \nabla_X\nabla_Y(fZ) - \nabla_Y\nabla_X(fZ) - \nabla_{[X,Y]}(fZ) \\
    &= \nabla_X\big((Yf)Z + f\nabla_Y Z\big) - \nabla_Y\big((Xf)Z + f\nabla_X Z\big)
    - ([X,Y]f)Z - f\nabla_{[X,Y]}Z.
\end{align*}
Expanding and simplifying: the terms involving $Xf$, $Yf$, $XYf$, $YXf$ cancel
(using $[X,Y]f = XYf - YXf$), leaving $fR(X,Y)Z$.
\end{proof}

\section{Coordinate expression}

In local coordinates, $R(\partial_i, \partial_j)\partial_k = R^l_{kij}\, \partial_l$,
with:
\[
    R^l_{kij} = \partial_i \Gamma^l_{jk} - \partial_j \Gamma^l_{ik}
    + \Gamma^l_{im}\Gamma^m_{jk} - \Gamma^l_{jm}\Gamma^m_{ik}.
\]

One also defines the \emph{lowered curvature tensor}:
\[
    R_{ijkl} = g_{im} R^m_{jkl},
\]
so that $R_{ijkl} = \inner{R(\partial_k, \partial_l)\partial_j}{\partial_i}$.

\begin{keyformulas}[title=\textbf{Components of the Riemann tensor}]
\[
    R^l_{kij} = \partial_i \Gamma^l_{jk} - \partial_j \Gamma^l_{ik}
    + \Gamma^l_{im}\Gamma^m_{jk} - \Gamma^l_{jm}\Gamma^m_{ik}.
\]
In normal coordinates at $p$: $R^l_{kij}(p) = \partial_i \Gamma^l_{jk}(p)
- \partial_j \Gamma^l_{ik}(p)$.
\end{keyformulas}

\section{Symmetries of the Riemann tensor}

\begin{theorem}[Fundamental symmetries]
\label{thm:symmetries}
The lowered curvature tensor $R_{ijkl}$ possesses the following symmetries:
\begin{enumerate}
    \item \textbf{Antisymmetry in $(k,l)$:} $R_{ijkl} = -R_{ijlk}$.
    \item \textbf{Antisymmetry in $(i,j)$:} $R_{ijkl} = -R_{jikl}$.
    \item \textbf{Pair symmetry:} $R_{ijkl} = R_{klij}$.
    \item \textbf{First Bianchi identity (algebraic):}
          $R_{ijkl} + R_{iklj} + R_{iljk} = 0$.
\end{enumerate}
In intrinsic notation:
\begin{enumerate}
    \item $R(X,Y) = -R(Y,X)$.
    \item $\inner{R(X,Y)Z}{W} = -\inner{R(X,Y)W}{Z}$.
    \item $\inner{R(X,Y)Z}{W} = \inner{R(Z,W)X}{Y}$.
    \item $R(X,Y)Z + R(Y,Z)X + R(Z,X)Y = 0$.
\end{enumerate}
\end{theorem}

\begin{proof}
(1) is immediate from the definition ($R(X,Y) = -R(Y,X)$).

(2) Follows from metric compatibility. Indeed:
\begin{align*}
    \inner{R(X,Y)Z}{Z} &= \inner{\nabla_X\nabla_Y Z - \nabla_Y\nabla_X Z}{Z} \\
    &= X\inner{\nabla_Y Z}{Z} - \inner{\nabla_Y Z}{\nabla_X Z}
    - Y\inner{\nabla_X Z}{Z} + \inner{\nabla_X Z}{\nabla_Y Z} \\
    &= \frac{1}{2}XY\inner{Z}{Z} - \frac{1}{2}YX\inner{Z}{Z}
    = \frac{1}{2}[X,Y]\inner{Z}{Z} = \inner{\nabla_{[X,Y]}Z}{Z}.
\end{align*}
Hence $\inner{R(X,Y)Z}{Z} = 0$, which by polarization gives (2).

(3) follows from (1), (2), and (4).

(4) First Bianchi identity: using torsion-freeness and the definition, expand the
three cyclic terms and verify the sum vanishes.
\end{proof}

\begin{proposition}[Number of independent components]
In dimension $n$, the number of independent components of the Riemann tensor is:
\[
    \frac{n^2(n^2-1)}{12}.
\]
For $n = 2$: $1$; for $n = 3$: $6$; for $n = 4$: $20$.
\end{proposition}

\section{Bianchi identities}

\begin{theorem}[Second Bianchi identity (differential)]
\[
    (\nabla_X R)(Y,Z) + (\nabla_Y R)(Z,X) + (\nabla_Z R)(X,Y) = 0.
\]
In coordinates: $\nabla_m R^l_{kij} + \nabla_i R^l_{kmj} + \nabla_j R^l_{kim} = 0$.
\end{theorem}

\begin{proof}
Work in normal coordinates at a point $p$ where $\Gamma^k_{ij}(p) = 0$. The
covariant derivative of $R$ reduces to the partial derivative, and the result
follows from a direct computation using commutativity of partial derivatives.
\end{proof}

\begin{remark}
The second Bianchi identity plays a fundamental role: when contracted, it yields
the divergence-freeness of the Einstein tensor
$G_{ij} = R_{ij} - \frac{1}{2}Sg_{ij}$ in general relativity.
\end{remark}

\section{Curvature and parallel transport}

\begin{theorem}[Interpretation via parallel transport]
Let $X, Y$ be two commuting vector fields ($[X,Y] = 0$). Consider the
infinitesimal parallelogram formed by the flows of $sX$ and $tY$. Parallel
transport of a vector $Z$ around this parallelogram yields, to second order:
\[
    Z_{\text{final}} - Z_{\text{initial}} = st\, R(X,Y)Z + O(s^2 + t^2).
\]
\end{theorem}

\begin{corollary}[Flatness]
$(M, g)$ is flat ($R \equiv 0$) if and only if parallel transport is
path-independent, if and only if $(M, g)$ is locally isometric to $\R^n$.
\end{corollary}

\section{Jacobi fields}

\begin{definition}[Jacobi field]
Let $\gamma$ be a geodesic. A vector field $J$ along $\gamma$ is a
\emph{Jacobi field} if it satisfies the \emph{Jacobi equation}:
\[
    \frac{D^2 J}{dt^2} + R(\dot\gamma, J)\dot\gamma = 0.
\]
\end{definition}

\begin{proposition}
Jacobi fields along a geodesic form a $2n$-dimensional vector space, determined
by the initial conditions $J(0)$ and $\frac{DJ}{dt}(0)$.
\end{proposition}

\begin{theorem}[Origin of Jacobi fields]
Jacobi fields are the variation fields of families of geodesics: if $\gamma_s$
is a smooth family of geodesics, then
$J(t) = \frac{\partial\gamma_s}{\partial s}\big|_{s=0}$ is a Jacobi field
along $\gamma_0$.
\end{theorem}

\begin{proof}
Set $T = \dot\gamma_s$ and $J = \frac{\partial\gamma_s}{\partial s}$.
Since $[T, J] = 0$ (symmetry of the variation) and $\nabla_T T = 0$ (geodesic
equation):
\[
    \frac{D^2 J}{dt^2} = \nabla_T\nabla_T J = \nabla_T\nabla_J T
    = \nabla_J\nabla_T T + R(T,J)T = R(T,J)T = -R(\dot\gamma, J)\dot\gamma.
\]
Hence $\frac{D^2 J}{dt^2} + R(\dot\gamma, J)\dot\gamma = 0$.
\end{proof}

\begin{example}[Jacobi fields on $S^n$]
On $S^n$ (curvature $K = 1$), along a unit-speed great circle $\gamma$, the
Jacobi equation for an orthogonal field $J \perp \dot\gamma$ becomes:
\[
    J''(t) + J(t) = 0,
\]
with solution $J(t) = A\cos t + B\sin t$. Conjugate points occur at
$t = k\pi$, $k \in \Z^*$.
\end{example}

\begin{example}[Jacobi fields on $\mathbb{H}^n$]
On $\mathbb{H}^n$ (curvature $K = -1$), for $J \perp \dot\gamma$:
\[
    J''(t) - J(t) = 0,
\]
with solution $J(t) = A\cosh t + B\sinh t$. There are \emph{no} conjugate points.
\end{example}

\section{Curvature computations on examples}

\begin{example}[Curvature of $S^2$]
In spherical coordinates:
\[
    R^1_{212} = \partial_1\Gamma^1_{22} - \partial_2\Gamma^1_{12}
    + \Gamma^1_{1m}\Gamma^m_{22} - \Gamma^1_{2m}\Gamma^m_{12}.
\]
One obtains $R^1_{212} = \sin^2\!\theta$, hence
$R_{1212} = g_{11} R^1_{212} = \sin^2\!\theta$. The Gaussian curvature is
$K = \frac{R_{1212}}{g_{11}g_{22} - g_{12}^2} = \frac{\sin^2\theta}{\sin^2\theta} = 1$.
\end{example}

\begin{example}[Curvature of $\mathbb{H}^2$]
With the metric $g = y^{-2}(dx^2 + dy^2)$, one computes $R_{1212} = -1/y^4$,
$\det(g) = 1/y^4$, giving $K = -1$.
\end{example}

\begin{example}[Curvature of a bi-invariant Lie group]
For left-invariant $X, Y, Z \in \mathfrak{g}$, with $\nabla_X Y = \frac{1}{2}[X,Y]$:
\[
    R(X,Y)Z = \frac{1}{4}[[X,Y],Z].
\]
\end{example}

\section{Second variation formula}

\begin{theorem}[Second variation of energy]
Let $\gamma : [0, L] \to M$ be a unit-speed geodesic and $\gamma_s$ a fixed-endpoint
variation with variation field $V$. Then:
\[
    \frac{d^2}{ds^2}\bigg|_{s=0} E(\gamma_s) = \int_0^L \left(
    \norm{\frac{DV}{dt}}^2 - \inner{R(\dot\gamma, V)\dot\gamma}{V}\right) dt.
\]
\end{theorem}

\begin{definition}[Index form]
The \emph{index form} of the geodesic $\gamma$ is the bilinear form:
\[
    I(V, W) = \int_0^L \left(\inner{\frac{DV}{dt}}{\frac{DW}{dt}}
    - \inner{R(\dot\gamma, V)\dot\gamma}{W}\right) dt.
\]
\end{definition}

\section{Exercises}

\begin{exercise}
Verify the first Bianchi identity $R(X,Y)Z + R(Y,Z)X + R(Z,X)Y = 0$ for $S^2$
in spherical coordinates.
\end{exercise}

\begin{exercise}
Show that in dimension $2$, the Riemann tensor is entirely determined by the
Gaussian curvature $K$:
$R_{1212} = K\,(g_{11}g_{22} - g_{12}^2) = K\det(g)$.
\end{exercise}

\begin{exercise}
Compute the Jacobi fields along a radial geodesic of the hyperbolic plane, and
verify that there are no conjugate points.
\end{exercise}

\begin{exercise}
For the Lie group $SO(3)$ with bi-invariant metric, compute $R(X,Y)Z$ for an
orthonormal basis $(X, Y, Z)$ of $\mathfrak{so}(3)$.
\end{exercise}

\begin{exercise}
Show that if $(M,g)$ is flat, then every Jacobi field has the form
$J(t) = A + tB$ for parallel vectors $A$ and $B$.
\end{exercise}

\begin{exercise}
Compute the Riemann tensor of the warped product $g = dt^2 + f(t)^2 g_{S^{n-1}}$
in terms of $f$ and its derivatives.
\end{exercise}
